\section{Reproducibility and the limits of computation}
\label{sec:verification}

The accompanying package contains the complete typeset article, a single
flattened \TeX\ source, the modular sources, eleven verification programs,
their recorded JSON results, and an integration guide. The proofs in the
article establish the uniform statements. The programs provide exact
finite certificates and independent numerical diagnostics at explicitly
specified orders.

\subsection{Exact finite certificates}

The order-two verifier constructs the matrix directly from the inherited
binomial formula, without defining it by the proposed factorization.
For every $b=2,\ldots,14$ it checks the determinant, the rank-one
factorization, the Pad\'e inverse-row contraction, the rational extraction
row, the shifted Hankel product, and the final factorial cancellation.
It also checks the Pad\'e coefficient identities through the corresponding
orders and reconstructs the ordinary-series examples at $b=2,3,4$.

The order-four verifier checks the symbolic perturbation contractions,
the resulting determinant polynomial, and its exact parity calculation.
It also constructs the original source matrices for $b=3,\ldots,14$,
including the exceptional smallest pair $(4,6)$. The general shifted-Hankel
verifier checks all 45 pairs $1\le a\le b\le9$, comparing the source
entries with the factorization, the determinant signs, the size-$a$
reduction, and the weighted Bernoulli sums. The exact ranges supplement
the uniform proofs of both new sections.

The order-six program constructs the normalized $3\times3$ matrix with
an indeterminate $b$. It proves the determinant normalization, constructs
its numerator in the integer polynomial ring, and verifies the modulo-four
specialization used in the nonvanishing proof. Direct matrices are also
checked at $b=4,\ldots,12$, and the order-eight counterexample to the
uniform valuation pattern is computed exactly. A separate short program
checks the explicit $(4,6)$ ordinary tail identity. These computations
are finite polynomial certificates, not interpolation from integer samples.

The Tornheim symbolic verifier multiplies the Gamma and zeta Taylor
series in exact rational arithmetic. It verifies both elementary
restrictions on $J$, the full cubic direction formula, and both printed
cancellation identities. A separate construction builds the linear
restriction matrices for homogeneous polynomial germs through degree ten
and checks the stated kernel dimensions. An additional independent script
derives the three-ray cancellation directly from raw series and the two
elementary constraints; it also detects a deliberately changed coefficient.

The differential verifier checks every partition monomial at every
weight through twelve. It verifies its exact decomposition into normal
form and derivative, the zero normal form of every derivative column,
the factor-degree counts, the adjoint concomitant, and the Gamma parity.
It constructs the full degree-six and degree-seven polynomial identities,
including the canonical endpoint constants. The Laurent principal part
of the retained sixth mixed integral is checked symbolically as well.

The checks over finite ranges confirm the formulas within those
ranges. They do not by themselves prove the uniform statements.
The identities certified with indeterminate $b$, including the
integer-polynomial certificate in Section~\ref{rs:section}, are exact
finite algebraic steps in the proofs. The uniform conclusions also
use the Pad\'e coefficient argument, the Bernoulli congruences,
holomorphic divisibility, and the terminating partition argument,
as applicable. No external proof assistant was used to formalize
the article's analytic proofs.

\subsection{Independent numerical diagnostics}

Table~\ref{tab:numerics} reports the main recorded comparisons. Every
entry is floating-point evidence, not an interval enclosure. The numerical
evaluators do not assume the identities they are checking.

\begin{table}[htbp]
\centering\small
\caption{Recorded numerical comparisons; residuals are diagnostic.}
\label{tab:numerics}
\begin{tabular}{@{}p{.33\textwidth}p{.36\textwidth}>{\raggedleft\arraybackslash}p{.24\textwidth}@{}}
\toprule
Identity & Independent evaluation & Absolute residual\\
\midrule
Ordinary $(2,4)$ tail sum & 100-digit tails; 100 initial terms and 24 asymptotic terms
 & $9.93\cdot10^{-70}$\\[5pt]
Ordinary $(2,6)$ tail sum & Same method, independently formed weight polynomial
 & $3.20\cdot10^{-65}$\\[5pt]
Ordinary $(2,8)$ tail sum & 100 digits; 120 initial terms and 32 asymptotic terms
 & $3.28\cdot10^{-72}$\\[5pt]
Three cubic Tornheim rays & Subtracted Mellin evaluation and odd-part extrapolation
 & $6.30\cdot10^{-33}$\\[5pt]
Canonical sixth digamma formula & Direct Laurent subtraction and separate height-one Mellin integral
 & $5.63\cdot10^{-51}$\\[5pt]
$(\psi')^2$ and $(\psi'')^2$ & Direct Laurent subtraction and ordinary quadrature
 & $\begin{gathered}1.23\cdot10^{-90},\\1.77\cdot10^{-88}\end{gathered}$\\
\bottomrule
\end{tabular}
\end{table}

For the harmonic sums, the exact polynomial subtraction removes every
nonnegative asymptotic power; the remaining inverse powers are summed
by ordinary Hurwitz zeta values to accelerate a convergent series.
For the independent $(2,8)$ check, increasing the precision and the
tail parameters changed the residual from approximately
$4.20\cdot10^{-52}$ to $3.28\cdot10^{-72}$.

The Tornheim evaluator integrates the Jonqui\`ere expansion below $x=1$
and sums an incomplete-Gamma tail above $x=1$. With 50-digit arithmetic,
series index 48, tail index 110, and seven odd-part steps starting at
$h=0.0008$, it gives
\[
 W_{1,2,3}'''(0)+4W_{2,3,1}'''(0)-5W_{3,1,2}'''(0)
 \approx-451.3535195684862560644674145675405.
\]
The code does not use the proposed directional formula, fit $\Omega$
or $\kappa$, or derive its values from the digamma moment.

For the sixth identity, direct finite-part quadrature gives
\begin{align*}
 \FP\int_0^1P_6(\psi(x))\,dx
  &\approx20.8226131542578224267340551051333,\\
 \widetilde M_6
  &\approx-3.40244508656918661747498017448045.
\end{align*}
The separate Mellin integral uses an explicit lower cutoff $10^{-60}$;
the omitted piece is the limiting numerical effect in the displayed
comparison. Its size was not promoted to a certified error bound. The
formula itself follows from exact polynomial identities, exact endpoints,
and the proved ordinary Gamma bridge.

\subsection{Build and integration procedure}

Run \src{python build.py} to regenerate the flattened source and PDF.
The build uses \src{pdflatex} and resolves cross-references in repeated
passes. Run \src{python verification/verify_all.py} for the exact
certificates. The optional flag \src{--numerical} also runs the numerical
diagnostics; the Mellin computation is the longest job.
Required Python packages and the versions used are recorded in the
package. The JSON files contain parameters and full recorded output,
so a future replay can distinguish a change in arithmetic, a change in
numerical settings, and a change in the mathematical claim.

The integration guide maps each contribution to its source question and
names the labels that may be imported into the manuscript. It proposes
status edits without changing the pinned source checkout. The article is
intended to be read and checked as a research continuation before those
edits are adopted.
